\documentclass[journal]{IEEEtran}



\ifCLASSINFOpdf

\else

\fi



\usepackage{inconsolata} 
\usepackage[most]{tcolorbox}
\usepackage{xcolor}
\usepackage{listings}
\usepackage{amssymb}
\usepackage{amsmath}
\usepackage[table]{xcolor}
\usepackage{url}
\usepackage{graphicx} % Added to allow \includegraphics to work

% 1. Define Colors
\definecolor{codebg}{rgb}{1,1,1}       % Pure White
\definecolor{framecolor}{rgb}{0.95,0.95,0.95} % Very light grey for the header
\definecolor{keywordblue}{rgb}{0.13,0.13,1}
\definecolor{typegreen}{rgb}{0,0.6,0}
\definecolor{commentgray}{rgb}{0.5,0.5,0.5}

% 2. Define ProVerif Language for Syntax Highlighting
\lstdefinelanguage{ProVerif}{
  keywords={let, in, out, if, then, else, new, event, fun, type, query, reduction, process, letfun, equation, not, compfun, forall, otherwise},
  keywordstyle=\color{keywordblue}\bfseries,
  ndkeywords={bitstring, t_sk, t_pk, t_exponent, G, t_key},
  ndkeywordstyle=\color{typegreen}\bfseries,
  identifierstyle=\color{black},
  sensitive=false,
  comment=[l]{//},
  morecomment=[s]{(*}{*)},
  commentstyle=\color{commentgray}\ttfamily,
  stringstyle=\color{purple}\ttfamily,
}

% 3. Create the "proverifcode" Environment
\newtcblisting[auto counter]{proverifcode}[2][]{
    enhanced,
    listing only,
    listing engine=listings,
    listing options={
        language=ProVerif,
        basicstyle=\ttfamily\scriptsize, 
        breaklines=true,
        tabsize=2,
    },
    title={#2}, 
    colback=codebg,     
    colframe=framecolor, 
    coltitle=black, 
    sharp corners, 
    boxrule=0.5pt,      
    left=2mm, right=2mm, top=2mm, bottom=2mm, 
    #1
}
% -------------------------------------------------------------



\hyphenation{Project Title}


\begin{document}

\title{Hash Gone Right: Completing ProVerif's Resolution for Hash-Based Protocols}


\author{Arun~Natarajan
        and~Hafeez~Muhammed
}



% The paper headers
\markboth{B.Tech CSE Semester 8 Mid-Term Evaluation Report}%
{}






% make the title area
\maketitle

% As a general rule, do not put math, special symbols or citations
% in the abstract or keywords.
\begin{abstract}
Symbolic verification tools like ProVerif traditionally model hash functions as ideal random oracles, ignoring the associative structure of real-world Merkle-Damg\aa{}rd constructions. Previous attempts to model this associativity, such as the ``Hash Gone Bad'' framework, rely on bounded recursive functions that suffer from combinatorial explosions and yield inconclusive verdicts for complex protocols like IKE. To resolve this, we integrate Jaffar’s Minimal and Complete Word Unification algorithm into ProVerif’s axiom system, replacing naive variable splitting with efficient directed splitting. Furthermore, because standard unification is insufficient for ProVerif's abstract resolution phase, we introduce two complementary sets of resolution-level axioms: MDH associativity merge axioms for depth mismatches and attacker-head axioms for attacker-controlled variables. This enhanced framework eliminates artificial depth bounds, averts state-space explosions, and enables ProVerif to achieve definitive security verdicts where previous models could not.
\end{abstract}

% Note that keywords are not normally used for peerreview papers.
\begin{IEEEkeywords}
ProVerif, Symbolic Verification, Cryptographic Protocols, Hash Functions, Merkle-Damg\aa{}rd, Unification.
\end{IEEEkeywords}



\IEEEpeerreviewmaketitle


\section{Introduction}

\IEEEPARstart{T}{he} security of cryptographic protocols relies heavily on hash functions. Verification tools like ProVerif traditionally model these as ideal random oracles, ignoring the associative structure of real-world Merkle-Damg\aa{}rd constructions where $H((a,b), c)$ is equivalent to $H(b, H(a, c))$. This abstraction masks vulnerabilities like chosen-prefix collisions. The ``Hash Gone Bad'' framework \cite{cheval2023hash} attempted to address this gap using bounded recursive computation functions. However, its manual variable decomposition causes a combinatorial explosion, and its artificial depth bounds lead to incomplete, ``cannot be proved'' verdicts for complex protocols like IKE.

To eliminate these bounds, we implemented Jaffar's Minimal and Complete Word Unification algorithm \cite{jaffar1990word} into ProVerif's axiom system. This approach replaces blind variable guessing with directed splitting, where hash equations are resolved through three principled cases of equality, left-prefix, and right-prefix overlap.

Yet, even with Jaffar's algorithm in place, several IKE properties still returned ``cannot be proved.'' To understand why, we must consider ProVerif's two-phase verification architecture. During Phase 1 (resolution), ProVerif builds an abstract over-approximation of reachable states. If it detects a potential attack here but fails to instantiate it into a concrete execution trace in Phase 2, the abstract attack is deemed a spurious artifact, resulting in the inconclusive ``cannot be proved'' verdict. The root problem, therefore, is not in trace reconstruction itself, but in the resolution phase being too imprecise: it admits abstract attack paths that do not correspond to any real execution.

We identified two specific sources of this imprecision in the context of Merkle-Damg\aa{}rd hash reasoning. First, the axiom system lacks merge axioms for depth-mismatch equations. When MDH normalization maps a shallow chain to a deeper one (e.g., $H(x, Nil) = H(y, H(z, Nil))$ requires $x = (z, y)$), no existing axiom captures this, leaving hash terms that should unify stuck as unresolved constraints. Second, the \texttt{instantiate\_allowed} termination guard, designed to prevent non-termination by blocking variable splitting when the opposing chain contains variables, is overly conservative when the attacker controls a head variable. In cross-session scenarios where the attacker supplies a hash block, the guard blocks the only axiom that could unify the attacker's term with the honest party's value, leaving the equation permanently unresolved.

This paper introduces two complementary sets of resolution-level axioms to close these gaps. First, MDH associativity merge axioms handle depth-mismatch cases by directly encoding algebraic identities such as $H((z,y), Nil) = H(y, H(z, Nil))$. Second, attacker-head axioms bypass the conservative termination guard to force head equality specifically when an attacker-controlled variable faces a non-variable term, without introducing non-termination. Together, these axioms sharpen the resolution phase enough for ProVerif to reach definitive verdicts on all previously inconclusive IKE properties. The remainder of this paper details ProVerif's architecture and relevant unification algorithms (Section II), analyzes the ``cannot be proved'' problem and its root causes (Section III), presents our axiom contributions (Section IV), discusses experimental results (Section V), and concludes with future work (Section VI).

\section{Literature Survey}

\subsection{Cryptographic Hashes and the Random Oracle Model}
The security of modern cryptographic protocols is critically dependent on the integrity of their underlying cryptographic primitives, particularly hash functions. Many analysis techniques rely on the Random Oracle Model (ROM), which assumes a perfect hash function. This idealization ignores the structural properties and iterative processing that characterize widely deployed algorithms such as MD5 and SHA-1.  This creates a significant gap between a protocol's proven security in an idealized model and its actual vulnerabilities to length-extension or collision attacks in real-world implementations. Real-world Merkle-Damg\aa{}rd constructions are inherently order-sensitive, processing blocks iteratively such that terms exhibit associative properties: $H((a,b), c)$ is algebraically equivalent to $H(b, H(a, c))$.

\subsection{The ``Hash Gone Bad'' Framework}
The foundational work in the ``Hash Gone Bad'' framework \cite{cheval2023hash} directly addressed this gap. The approach extended ProVerif with recursive computation functions to approximate the associative nature of Merkle-Damg\aa{}rd constructions. This solution relies on bounded approximations and manual axioms that limit the depth of hash analysis to prevent infinite recursion. The logic follows a flattening strategy implemented via recursive pattern matching.

% Time Complexity: O(3^d) verification time due to combinatorial branching up to depth d.
% Space Complexity: Exponential memory overhead during state exploration.
\begin{proverifcode}{Recursive MDH Logic}
compfun MDH(bitstring):bitstring =
  forall x:bitstring; MDH(x) if is_var(x) || x = Nil -> x
  otherwise forall h1,h2,h3:bitstring; 
    MDH(H(CPcol1(h1,h2),h3)) -> H(CPcol1(MDH(h1),MDH(h2)),MDH(h3))
  otherwise forall x1,x2,h:bitstring; 
    MDH(H((x1,x2),h)) -> MDH(H(x2,H(x1,h))) (* Associativity Rule *)
  otherwise forall x,h:bitstring; 
    MDH(H(x,h)) -> H(x,MDH(h))
  [mayFail]
.
\end{proverifcode}

While this bounded associativity allows for some analysis, it creates a severe trade-off between completeness and computational efficiency. Enabling collisions produces dramatic slowdowns; for example, Simplified IKE verification time jumps significantly compared to the baseline. This combinatorial explosion stems from manual axioms that blindly decompose variables (e.g., guessing $x = (y,z)$), forcing the solver to explore an exponentially large space of invalid hash structures. Furthermore, once a protocol requires reasoning beyond the preset manual depth bounds, ProVerif aborts the reasoning process and returns an inconclusive ``cannot be proved'' outcome.

\subsection{Unification Algorithms for Associativity}
To address the limitations of bounded associativity, three distinct classes of unification algorithms must be considered. Associative-Commutative (AC) unification solves equations where operators satisfy both associativity and commutativity. However, cryptographic hash functions are inherently order-sensitive ($H(a,b) \neq H(b,a)$), making the assumption of commutativity unsound. Plotkin's recursive descent algorithm provides a naive approach to strictly associative unification. While algebraically sound, it is infinitary and produces infinitely many redundant unifiers without a termination guarantee.

Jaffar's algorithm for Minimal and Complete Word Unification \cite{jaffar1990word} introduces a significant refinement. Instead of blind recursive splitting, it employs directed splitting.  Given an equation $x \cdot u = y \cdot v$, the algorithm proves that solutions exist only if one variable is a prefix of the other. This restricts the search space to three mathematically necessary cases: equality ($x=y$), left-prefix overlap ($y = x \cdot z$), and right-prefix overlap ($x = y \cdot z$). This approach avoids state-space explosion and guarantees termination without arbitrary bounds. To prevent non-termination during resolution, the axiom system employs a termination guard, \texttt{instantiate\_allowed}, which blocks variable instantiation when the opposing hash chain contains unresolved variables.

% Time Complexity: O(1) per split step, worst-case linear O(n) relative to chain depth.
% Space Complexity: O(n) for the unification call stack.
\begin{proverifcode}{Adapted Word Unification for Hashes}
(* Input: Equation H(x, h1) = H(y, h2) *)
(* Output: Set of unifiers {sigma} *)

If x and y are identical:
  Return { x -> x } AND unify(h1, h2)
Else:
  Attempt Split 1 (x is prefix):
    Let y = (x, z) where z is a fresh variable
    Return { y -> (x, z) } AND unify(h1, (z, h2))
  Attempt Split 2 (y is prefix):
    Let x = (y, z) where z is a fresh variable
    Return { x -> (y, z) } AND unify((z, h1), h2)
\end{proverifcode}

We implemented Jaffar's algorithm into ProVerif's axiom system, resolving the performance bottleneck. However, as we show in Section III, the resulting axiom system still contains gaps that produce inconclusive verdicts on IKE.

\subsection{ProVerif's Verification Architecture}
To understand why these inconclusive outcomes occur, it is necessary to examine ProVerif's two-phase verification architecture.  In the first phase, known as the resolution phase, ProVerif saturates Horn clauses to build an abstract over-approximation of all reachable protocol states. If no potential attack is derived during this phase, the security property is proved true. If a potential attack is found, ProVerif enters the second phase, trace reconstruction. In this phase, the tool attempts to instantiate the abstract derivation into a concrete, executable attack trace. If reconstruction succeeds, a real attack exists. If it fails, the abstract attack was a spurious artifact of the over-approximation, and ProVerif reports the inconclusive ``cannot be proved'' verdict.

\begin{table*}[t]
\centering
\caption{IKEv2 Query Results with Jaffar's Axioms}
\label{tab:query_results}
\renewcommand{\arraystretch}{1.3} 
\begin{tabular}{|l|l|l|l|}
\hline
\textbf{Query} & \textbf{Description} & \textbf{Assoc, No Col} & \textbf{Assoc + Col} \\ \hline
Q1 & Type-flaw agreement & false (attack) & false (attack) \\ \hline
Q2 & Transcript agreement & cannot be proved & cannot be proved \\ \hline
Q3 & Hash agreement & cannot be proved & cannot be proved \\ \hline
Sanity1 & Initiator reachable (cookies forbidden) & false $\checkmark$ & false $\checkmark$ \\ \hline
Sanity2 & Initiator reachable (cookies allowed) & false $\checkmark$ & false $\checkmark$ \\ \hline
Sanity3 & Responder reachable (cookies forbidden) & false $\checkmark$ & false $\checkmark$ \\ \hline
Sanity4 & Responder reachable (cookies allowed) & false $\checkmark$ & false $\checkmark$ \\ \hline
\end{tabular}
\end{table*}

\section{Diagnosing the Remaining Imprecisions}

After integrating Jaffar's directed-splitting axioms into the \texttt{hash\_no\_collision.pvl} and \texttt{hash\_collision.pvl} libraries, we evaluated the framework on the Simplified IKEv2 protocol under both the no-collision and chosen-prefix collision threat models. The protocol specification defines three security queries of increasing strength:

\begin{itemize}
    \item \textbf{Q1 (Type-Flaw Agreement):} Verifies that if responder $B$ accepts a signed payload $s_A$ from initiator $A$, then $A$ did indeed produce $s_A$. This query captures \emph{type-flaw attacks}, where the attacker reinterprets the bit-level encoding of a message to forge a valid signature. 
    \item \textbf{Q2 (Transcript Agreement):} Verifies that the transcript hashes computed independently by $A$ and $B$ satisfy \texttt{eq\_hash}$(t, t')$. This ensures both parties agree on the \emph{input} to the hash function.
    \item \textbf{Q3 (Hash Agreement):} Verifies the stronger property that \texttt{eq\_hash}$(h, h')$, where $h$ and $h'$ are the final hash outputs. Since normalization via \texttt{MDH} ensures $t = h$ when the inputs are identical, Q2 and Q3 are logically equivalent in the no-collision model but may diverge under collisions.
\end{itemize}

In addition, four \emph{sanity queries} confirm that each role (initiator and responder) is reachable under both cookie-enabled and cookie-disabled modes. The empirical results of these queries under Jaffar's directed splitting are summarized in Table \ref{tab:query_results}.

With Jaffar's axioms in place, Q1 correctly returns \textbf{false} (attack found), confirming the known type-flaw vulnerability. However, Q2 and Q3 both return \textbf{``cannot be proved''}, an inconclusive result indicating that ProVerif's resolution engine was unable to either confirm or refute the property. All four sanity queries return false as expected, ruling out a modeling error. We traced this failure to two distinct imprecisions in the axiom system, both operating at the resolution level (Phase 1).

\subsection{Root Cause 1: Depth Mismatch in Merge Resolution}

Inspection of ProVerif's internal abstract goals revealed that the resolution engine generates subgoals of the form:

$$\texttt{eq\_hash}\bigl(H(x,\, \texttt{Nil}),\; H(y,\, H(z,\, \texttt{Nil}))\bigr)$$

where the left-hand chain has depth 1 and the right-hand chain has depth 2.  Under MDH normalization, these two terms \emph{are} equal when $x = (z, y)$, because:

$$H\bigl((z,\,y),\, \texttt{Nil}\bigr) \;\xrightarrow{\text{MDH}}\; H\bigl(y,\, H(z,\, \texttt{Nil})\bigr)$$

However, Jaffar's existing axioms only fire when both sides of the equation have matching structural patterns (variable-headed chains of equal depth). Because the depths differ, none of the directed-splitting rules apply, and the goal remains unresolved, as visualized in Fig. \ref{fig:root_1}.

\begin{figure}[!ht]
    \centering
    \includegraphics[width=1\linewidth]{Root_1.png}
    \caption{Visual representation of the Depth Mismatch problem in MDH chain unification during the abstract resolution phase.}
    \label{fig:root_1}
\end{figure}

\subsection{Root Cause 2: The \texttt{instantiate\_allowed} Guard}

Even after accounting for depth mismatches, a second class of goals remained stuck. (The iterative construction of the transcript hash, which generates these complex structures, is depicted in Fig. \ref{fig:ike_overview}.) ProVerif generates abstract goals such as:

$$\texttt{eq\_hash}\bigl(H(\mathit{sA}_5,\, H(\mathit{ck}_2,\, \texttt{Nil})),\; H(\mathit{sA}_4,\, H(\mathit{cookie}_1,\, \texttt{Nil}))\bigr)$$

Here $\mathit{sA}_5$ is an \emph{attacker-controlled variable} (the attacker may choose its value), $\mathit{sA}_4$ is a protocol name (not a variable), $\mathit{ck}_2$ is the initiator's cookie, and $\mathit{cookie}_1$ is the responder's cookie variable. The natural resolution is to deduce $\mathit{sA}_5 = \mathit{sA}_4$, which would reduce the goal to $\texttt{eq\_hash}(H(\mathit{ck}_2, \texttt{Nil}),\, H(\mathit{cookie}_1, \texttt{Nil}))$.

This deduction is blocked by the \texttt{instantiate\_allowed} guard in the existing axiom system.  The guard function traverses the opposing chain and returns \texttt{false} whenever it encounters a variable in a head position, which in this case is $\mathit{cookie}_1$. This conservative check was introduced to prevent non-termination: without it, the axiom $x_1 = x_2 \,\|\, x_1 = (x_1', x_1'')$ could trigger an infinite chain of fresh variable introductions when both sides contain unresolved variables. However, the guard is \emph{overly conservative} when one of the head variables is attacker-controlled, because the attacker's variable can be safely instantiated to a concrete value without risk of divergence. This blockage is visually demonstrated in Fig. \ref{fig:root_2}.

\begin{figure*}[t] % Added [t] to suggest top-of-page placement
    \centering
    % Change 1\linewidth to a decimal like 0.5\textwidth
    \includegraphics[width=0.85\textwidth]{hash_construction.png.png}
    \caption{Overview of the Simplified IKEv2 protocol execution and the iterative construction of the transcript hash.}
    \label{fig:ike_overview}
\end{figure*}

\begin{figure}[!ht]
    \centering
    \includegraphics[width=1\linewidth]{root_2.png}
    \caption{Abstract resolution goal blocked by the \texttt{instantiate\_allowed} termination guard when facing attacker-controlled variables.}
    \label{fig:root_2}
\end{figure}

\section{Proposed Solution: Resolution-Level Axioms}

The two imprecisions diagnosed in Section III both occur during ProVerif's resolution phase (Phase 1) before trace reconstruction begins. Therefore, the fix must take the form of resolution-level axioms, which are Horn clause rules that fire during clause saturation to guide the engine toward the correct unifier. Two complementary families of axioms are introduced: \emph{merge axioms} to close the depth-mismatch gap, and \emph{attacker-head axioms} to bypass the overly conservative \texttt{instantiate\_allowed} guard.

\subsection{MDH Associativity Merge Axioms}

The core property exploited by the merge axioms is the MDH normalization rule:
\[
  H\bigl((a, b),\; L\bigr) \;\xrightarrow{\text{MDH}}\; H\bigl(b,\; H(a,\; L)\bigr)
\]
When ProVerif generates a goal such as $\texttt{eq\_hash}\bigl(H(x,\, \texttt{Nil}),\; H(y,\, H(z,\, \texttt{Nil}))\bigr)$, the axiom must ``teach'' the engine that the depth-1 side can represent a packed tuple that, once normalized, matches the depth-2 side. Applying the normalization rule in reverse:
\[
  H\bigl(y,\, H(z,\, \texttt{Nil})\bigr) \;\xleftarrow{\text{MDH}}\; H\bigl((z, y),\, \texttt{Nil}\bigr)
\]
yields the substitution $x = (z, y)$.  This logic is encoded as a pair of symmetric axioms (one for each argument order):

% Time Complexity: O(1) per application during the resolution phase.
% Space Complexity: O(1) memory footprint per unification step.
\begin{proverifcode}{Depth-1 vs Depth-2 Merge Axiom}
axiom x, y, z: bitstring;
  eq_hash(H(x, Nil), H(y, H(z, Nil))) ==>
    x = (z, y)
.
axiom x, y, z: bitstring;
  eq_hash(H(y, H(z, Nil)), H(x, Nil)) ==>
    x = (z, y)
.
\end{proverifcode}

The same reasoning extends to a depth-3 chain. Applying MDH normalization twice:
\begin{align*}
  H\bigl(y,\, H(z,\, H(w,\, \texttt{Nil}))\bigr) 
  &\xleftarrow{\text{MDH}} H\bigl((z,y),\, H(w,\, \texttt{Nil})\bigr) \\
  &\xleftarrow{\text{MDH}} H\bigl((w,(z,y)),\, \texttt{Nil}\bigr)
\end{align*}
yields the substitution $x = (w, (z, y))$:

\begin{proverifcode}{Depth-1 vs Depth-3 Merge Axiom}
axiom x, y, z, w: bitstring;
  eq_hash(H(x, Nil), H(y, H(z, H(w, Nil)))) ==>
    x = (w, (z, y))
.
axiom x, y, z, w: bitstring;
  eq_hash(H(y, H(z, H(w, Nil))), H(x, Nil)) ==>
    x = (w, (z, y))
.
\end{proverifcode}

\begin{table*}[h!]
\centering
\caption{Simplified IKEv2: Before vs. After Axiom Integration}
\label{tab:ike-results}
\renewcommand{\arraystretch}{1.3} 
\begin{tabular}{|l|c|c|c|c|}
\hline
 & \multicolumn{2}{c|}{\textbf{Baseline (Before)}} & \multicolumn{2}{c|}{\textbf{With Axioms (After)}} \\ \cline{2-5}
\textbf{Query} & \textbf{No Col} & \textbf{Col} & \textbf{No Col} & \textbf{Col} \\ \hline
Q1: Type-flaw      & false & false & false & false \\ \hline
Q2: Transcript     & cannot be proved & \cellcolor{red!15}cannot be proved & \cellcolor{green!15}\textbf{true} & \cellcolor{green!15}\textbf{true} \\ \hline
Q3: Hash           & cannot be proved & \cellcolor{red!15}cannot be proved & \cellcolor{green!15}\textbf{true} & \cellcolor{green!15}\textbf{true} \\ \hline
Sanity 1--4        & false & false & false & false \\ \hline
\end{tabular}
\end{table*}

Identical axioms are added for the \texttt{eq\_hash\_col} predicate in the collision library. In practice, depth-1 vs depth-2 and depth-1 vs depth-3 cover all the abstract goals generated by the Simplified IKE and SIGMA protocol models. Deeper mismatches would require additional axioms following the same pattern, but none were observed during our benchmark evaluations.

\subsection{Attacker-Head Axioms}

The second axiom addresses the case where a head variable is attacker-controlled and the opposing head is a concrete protocol term. As shown in Section III-B, the existing \texttt{instantiate\_allowed} guard rejects this case because a \emph{different} variable (e.g., a cookie) appears elsewhere in the opposing chain. The key insight is that an attacker-controlled variable can be safely instantiated to any concrete value; the attacker simply ``chooses'' the right value, so the guard's non-termination concern does not apply.

The axiom checks three conditions: (1) the head is a variable (\texttt{is\_var}), (2) the attacker controls it (\texttt{attacker}), and (3) the opposing head is not a variable (\texttt{not(is\_var)}). When all three hold, the variable is unified with the opposing head:

% Time Complexity: O(1) per evaluation during Phase 1.
% Space Complexity: O(1) memory footprint.
\begin{proverifcode}{Attacker-Head Resolution Axiom}
axiom x1, h1, x2, h2: bitstring;
  eq_hash(H(x1,h1),H(x2,h2))
    && is_var(x1) && attacker(x1) && not(is_var(x2))
    ==> x1 = x2;
  eq_hash(H(x1,h1),H(x2,h2))
    && is_var(x2) && attacker(x2) && not(is_var(x1))
    ==> x2 = x1
.
\end{proverifcode}

\textbf{Design choice: Omitting the pair branch.} The original Jaffar axioms include a disjunctive branch $x_1 = x_2 \;\|\; x_1 = (x_1', x_1'')$ to account for depth mismatches. In the attacker-head axiom, this pair branch was initially included but caused \emph{non-termination}: starting from $x_1 = (x_1', x_1'')$, the resolution engine would generate a fresh variable $x_1'$ which could again be attacker-controlled, triggering the same axiom indefinitely. Removing the pair branch is safe because any depth mismatch is already handled by the merge axioms detailed in Section IV-A. The two axiom families are therefore complementary: merge axioms resolve structural differences, and the attacker-head axiom resolves variable-instantiation blockages.

\textbf{Collision model adaptation.} The same attacker-head axiom is replicated for the \texttt{eq\_hash\_col} predicate in \texttt{hash\_collision.pvl}.


\section{Experimental Results and Evaluation}

To rigorously evaluate the effectiveness of the proposed resolution-level axioms, we analyzed the Simplified IKEv2 protocol under two distinct threat models: \emph{no collision} (associativity only) and \emph{chosen-prefix collision}. All experiments were executed within an isolated Docker environment  \cite{our_docker} to ensure reproducibility. The baseline configuration corresponds to the original Cheval et al. framework \cite{cheval2023hash} integrated with Jaffar's directed-splitting algorithm \cite{jaffar1990word}, but \emph{without} our newly introduced merge and attacker-head axioms.

Table \ref{tab:ike-results} presents a before-and-after comparative analysis for all seven Simplified IKEv2 security and reachability queries.


The empirical results demonstrate three critical outcomes of our A-Unification-enhanced framework:

\begin{itemize}
    \item \textbf{Definitive Security Verdicts (Q2 and Q3):} Both the transcript agreement and hash agreement queries successfully transition from the inconclusive ``cannot be proved'' state to a definitive \textbf{true}. This confirms that the IKE protocol is mathematically secure against transcript manipulation under both the no-collision and chosen-prefix collision threat models. Eliminating these resolution-level bottlenecks represents the primary contribution of our axiom sets.

    \item \textbf{Preservation of Known Vulnerabilities (Q1):} The framework correctly preserves the detection of the known type-flaw attack (returning \textbf{false}). This confirms that our axioms do not introduce unsound algebraic constraints that might mask legitimate protocol flaws. The vulnerability arises from the protocol's flat-tuple structural encoding and remains strictly independent of the hash unification logic.

    \item \textbf{Soundness and Reachability (Sanity 1--4):} All four sanity queries return false (indicating that the target states remain reachable). This verifies that the new axioms successfully sharpen the resolution phase without inadvertently restricting valid execution paths or causing false-positive verification halts.
\end{itemize}


% Time Complexity: O(1) for compilation of this static LaTeX text.
% Space Complexity: O(1) memory footprint for rendering.
\section{Conclusion and Future Work}

This work addresses a fundamental limitation in the symbolic verification of Merkle-Damg\aa{}rd hash constructions. While previous frameworks relied on bounded approximations that produced inconclusive results, our integration of Jaffar's minimal and complete word unification algorithm \cite{jaffar1990word} replaces blind variable splitting with principled \textbf{directed splitting}. 



By closing the \textbf{depth-mismatch gap} and bypassing the conservative \textbf{\texttt{instantiate\_allowed} guard} with targeted resolution-level axioms, we have restored both completeness and tractability to the analysis. Our proposed \textbf{merge axioms} and \textbf{attacker-head axioms} successfully resolved previously inconclusive queries for the Simplified IKEv2 protocol, transitioning \textbf{Q2 (transcript agreement)} and \textbf{Q3 (hash agreement)} to definitive \textbf{true} verdicts. Furthermore, \textbf{no regressions} were observed, as known type-flaw vulnerabilities and sanity reachability remained correctly identified.

\vspace{0.3cm} 

Future work will focus on \textbf{broadening the validation} of these axioms against a more diverse class of protocols to ensure universal coverage. 



Specifically, we aim to evaluate the framework against transcripts of higher arity and complex interleaved hash comparisons. This will help determine if further refinements, such as automated generation of deeper merge levels, can be integrated to establish a fully \textbf{autonomous and complete framework for algebraically-aware protocol verification}.


% needed in second column of first page if using \IEEEpubid
%\IEEEpubidadjcol










\section*{Acknowledgment}

We would like to express our sincere gratitude to our project advisor, Dr. Vinod Pathari, for his invaluable guidance, continuous support, and insightful feedback throughout the course of this research. His expertise was instrumental in shaping the direction of this project.

We also wish to thank the authors of ``Hash Gone Bad,'' \cite{cheval2023hash} whose foundational work provided the primary inspiration and basis for our research. Finally, we are grateful to the Department of Computer Science \& Engineering at NIT Calicut for providing us with the resources and academic environment necessary to complete this work.


\ifCLASSOPTIONcaptionsoff
  \newpage
\fi



\begin{thebibliography}{1}

\bibitem{cheval2023hash}
V.~Cheval, C.~Cremers, A.~Dax, L.~Hirschi, C.~Jacomme, and S.~Kremer, ``Hash Gone Bad: Automated Discovery of Protocol Attacks that Exploit Hash Function Weaknesses,'' in \emph{Proc. 32nd USENIX Security Symposium}, 2023.

\bibitem{jaffar1990word}
J.~Jaffar, ``Minimal and Complete Word Unification,'' \emph{IBM Thomas J. Watson Research Center}, Research Report RC 15663, 1990. [Note: Often published in Theoretical Computer Science, but citing the IBM report is standard for this foundational work.]

\bibitem{orig_repo}
V.~Cheval \emph{et al.}, ``Docker image and models for 'Hash Gone Bad','' GitHub Repository, 2023. [Online]. Available: \url{https://github.com/charlie-j/symbolic-hash-models}

\bibitem{our_repo}
A.~Natarajan and H.~Muhammed, ``Hash Gone Good,'' GitHub Repository, 2025. [Online]. Available: \url{https://github.com/Hafeez-hm/Fast_and_FuriHash}

\bibitem{our_docker}
A.~Natarajan, ``ProVerif with Computation Functions,'' Docker Hub Repository, 2025. [Online]. Available: \url{https://hub.docker.com/repository/docker/arunnats2004/proverif-compfun/general}

\end{thebibliography}


% that's all folks
\end{document}